\section{Curvature-aware local work sampling}\label{app:curvature-work}
The force escort translates every local candidate by the same vector. It does
not change the shape of the proposal cloud. Molecular energy can vary more
sharply in some directions, such as bond stretching, than in softer collective
motions. We test whether estimating this local curvature produces better
importance weights at a fixed candidate-query budget, without changing the
target distribution.

\paragraph{Fitting local curvature.}
Flatten the anchor into $3N$ coordinates. Let the columns of
$Q_N\in\mathbb R^{3N\times d}$ form an orthonormal basis for the centered
space, with $d=3(N-1)$. A displacement $u\in\mathbb R^d$ is expressed in
angstroms in this basis, and $Q_Nu$ gives the corresponding atom-coordinate
displacement. Let $f_A$ be the anchor force in the same basis.

We draw eight independent pilot candidates. Their displacements form columns
of $D$, and the negative changes in force form columns of $G$. We approximate
$G\approx B_A D$ by fitting a symmetric matrix:
\begin{equation}
 \|BD-G\|_F^2+\lambda\|B-\kappa I\|_F^2,
 \qquad
 \kappa=\max\!\left(0,\frac{\langle D,G\rangle_F}{\|D\|_F^2}\right),
 \qquad \lambda=0.01\,\|D\|_F^2/K.
\end{equation}
Here $K$ counts pilot columns that satisfy the graph restriction. If no pilot
survives, we set $B_A=0$ and recover the original force translation. For a nonempty
pilot we floor $\|D\|_F^2$ at $10^{-20}$ in the scalar and ridge calculations;
exactly zero displacements also give zero fitted curvature. The scalar
$\kappa$ estimates an average nonnegative stiffness. The second term shrinks
poorly determined directions toward this isotropic estimate. We replace negative
eigenvalues of the fitted matrix by zero to obtain a positive semidefinite
$B_A$. The isotropic control uses $B_A=\kappa I$ directly. Pilot information
and the fitted map are fixed before drawing production candidates.

\paragraph{Changing the proposal mean and covariance.}
The fitted quadratic energy is
$E_+(A+Q_Nu)-E_+(A)\approx-f_A^\top u+u^\top B_Au/2$.
Multiplying its Boltzmann factor by the original Gaussian reference gives another
Gaussian with a shifted mean and adjusted covariance. Define
$P_A=I+\sigma_A^2B_A/(k_B T_{\rm phys})$, a dimensionless positive-definite
precision matrix. We map the source displacement $u$ to
\begin{equation}
 v=P_A^{-1/2}u+b_A,\qquad
 b_A=P_A^{-1}\frac{\sigma_A^2f_A}{k_B T_{\rm phys}},\qquad
 \log J_A=-\tfrac12\log\det P_A.
\end{equation}
The factor $P_A^{-1/2}$ contracts perturbations more in stiff directions, while
$b_A$ adjusts the force-guided mean shift for curvature. The Jacobian $J_A$
measures volume change in the $d$ independent coordinates. We keep the original
reference width and temperature, so only the proposal changes, not the restrained
target. With $X=A+Q_Nu$ and $Y=A+Q_Nv$, complete work is
\begin{equation}
 W_A=E_+(Y)-E_+(A)+k_B T_{\rm phys}
 [\log q_A(X)-\log q_A(Y)-\log J_A].
\end{equation}
Candidates outside the original same-graph support receive zero weight. This is
a local use of targeted and escorted work reweighting
\citep{jarzynski2002targeted,escorted2008}.

\paragraph{Why curvature can improve the weights.}
Write the actual energy as the fitted quadratic plus a remainder $R_A(v)$.
Substitution gives $W_A=-k_B T_{\rm phys}\log Z_{\rm quad}+R_A(v)$, where
$Z_{\rm quad}$ is the integral of the Gaussian reference times the quadratic
Boltzmann factor. The fitted quadratic cancels from the varying part of the
work; only the approximation error changes across allowed candidates. If energy
is exactly that quadratic, all allowed candidates have equal weights. Tests
verify this constant-work identity for an anisotropic quadratic energy and check
rotation equivariance of the fitted map.

\paragraph{Pilot comparison at equal query budgets.}
For each of eight existing training compositions in each run, we take the first
four eligible anchors without ranking by energy or ESS. The translation baseline
uses 24 fresh production candidates per anchor. Each curvature-fitted map uses
at most eight pilot energy/force queries and 16 fresh production candidates.
The maximum particle-query budget is therefore the same. Common anchor queries
and reflection evaluations are counted separately. We evaluate all production
energies, including those of graph-invalid candidates, and do not reuse pilot
samples as production samples for the fitted maps.

\begin{table}[h]
\centering
\caption{Local ESS with complete-work weights at a maximum of 24 particle
queries per anchor. Each run uses 32 anchors. Curvature methods reserve eight
queries for fitting and use 16 production candidates; translation uses all 24
for production.}
\begin{tabular}{lrrr}
\toprule
Escort & Production particles & ESS, run 1 & ESS, run 2\\
\midrule
Force translation &24&2.22&1.86\\
Isotropic curvature &16&3.55&3.94\\
Directional curvature &16&6.69&6.70\\
\bottomrule
\end{tabular}
\end{table}

Directional curvature improves ESS over translation on 61 of 64 anchors.
The mean increase is 4.66 effective samples, with composition-bootstrap 95\%
interval $[3.81,5.46]$. Same-graph retention is 99.51\%, compared with 99.54\%
for translation. Independent multivariate-Gaussian density calculations reproduce
the stored complete-work weights. We then construct full teacher pools for all
349 original eligible anchors using fresh production draws. Translation versus
curvature gives mean ESS of 2.32 versus 6.71 in run 1 and 2.33 versus 6.37 in
run 2. These results concern the local teacher; the next experiment tests whether
they improve the trained generator.

\paragraph{Does the teacher improvement transfer to generation?}
We train three physical students: uniform sampling of 24 force-translated
candidates, complete-work weighting of those same 24 candidates, and complete-work
weighting of 16 curvature-adjusted candidates. Each uses the original
1,000-update schedule and reference/anchor choices. We reuse the matched replay
model and add the physical-minus-replay parameter difference with coefficient
one. All 349 original anchors remain eligible. The two work-weighted methods
have equal maximum particle-query budgets; uniform force targets can omit
candidate-energy queries. This follow-up changes the candidate-pool size and
curvature fit. The main experiments continue to use the original, preselected
eight-candidate force model.

\begin{table}[h]
\centering
\caption{Generated outputs after training with the three teachers, evaluated on
24 compositions with 1,536 attempts per method across both runs. Energy changes
compare a work-weighted student with the uniform-force student on paired, jointly
valid outputs. All methods use 128 network evaluations and no geometry optimization.}
\small
\begin{tabular}{lrrr}
\toprule
Teacher & Graph validity (\%) & eSEN change & GFN2 change\\
 & & (eV/atom) & (eV/atom)\\
\midrule
Uniform force, 24 candidates &54.30&0&0\\
Complete work, 24 candidates &53.39&$+0.00959$&$+0.00845$\\
Curvature work, 16 candidates &53.32&$+0.00793$&$+0.00783$\\
\bottomrule
\end{tabular}
\end{table}

Curvature work minus translation work changes energy by $-0.00184$ eV/atom under
eSEN (95\% interval $[-0.00654,0.00289]$) and $-0.00068$ under GFN2
($[-0.00484,0.00341]$). The two runs have opposite signs. The higher teacher
ESS therefore does not establish an additional generation benefit. Uniform force
targets give lower energy and force RMS than either work-weighted student under
both potentials. Relative to the frozen harmonic generator, this 24-candidate
uniform-force model lowers energy by $0.02088$ eV/atom under eSEN and $0.01918$
under GFN2. The primary method retains the inexpensive force-only recipe.
